% AI-assisted (Claude, writer), 26 Nov 2026: restructure after R1-ALIGN-OUTCOME (task writer/20261126-01).
% Text marked MOVED is carried over unchanged from the 23 Nov main text unless a comment says otherwise.
% Not an author decision. No label, test, evaluation or audit output was opened to write it.
\section{Full Reports Moved from the Main Text}
\label{app:results-main}

The main text reports the closed loop first and summarises the offline gate and the control errors. This appendix holds their full reports as drafted for both outcomes of the gate. Every value is a placeholder until its evidence is released.

% MOVED from results.tex (C1). One placeholder is re-worded for R1-STEP-ORDER (2)(d).
\subsection{Control errors of the default loop (C1)}
\label{app:c1}

Both sweeps ran the default loop: every proposal was executed once and kept if its observed metric beat the incumbent's recorded value. We score these logged decisions against the fresh-seed audits of Section~\ref{sec:setup}, on the eight T0 nodes and for both proposers. \TBD{Per the signed amendment (R1-POSTFREEZE-READS-CHOICES) as R1-STEP-ORDER (2)(d) reads it: all audited T0 instances of both sweeps with all ten seeds, read after the T0 test release row; if that row does not exist when the submission version is assembled, train and dev instances with the label seeds only, marked partial.} Table~\ref{tab:c1} gives the four rates of Section~\ref{sec:problem} per node and sweep. The two sweeps differ in the proposer and in its serving stack (Section~\ref{sec:setup}), so a difference between the sweep columns is a difference of both.

\begin{table*}[t]
\caption{Control errors of the default loop per T0 node, scored against fresh-seed audits, for sweep~1 (S1, the 7B manager) and sweep~2 (S2, the main proposer). Reading guide: a false keep is a kept candidate whose audited increment is not positive; a missed keep is a discarded candidate whose audited increment is positive; record bias is the recorded gain of the incumbent minus its audited gain, in units of $\sigma_{\mathrm{node}}$; the wasted-run share counts runs on no-ops, duplicates and invalid proposals. A difference between S1 and S2 is a difference of proposer and serving stack. \TBD{Intervals by instance-cluster bootstrap; node set of the pooled row; layout to be confirmed.}}
\label{tab:c1}
\centering
\small
\setlength{\tabcolsep}{4pt}
\begin{tabular}{lcccccccc}
\toprule
 & \multicolumn{2}{c}{False-keep rate} & \multicolumn{2}{c}{Missed-keep rate} & \multicolumn{2}{c}{Record bias} & \multicolumn{2}{c}{Wasted-run share} \\
\cmidrule(lr){2-3}\cmidrule(lr){4-5}\cmidrule(lr){6-7}\cmidrule(lr){8-9}
Node & S1 & S2 & S1 & S2 & S1 & S2 & S1 & S2 \\
\midrule
digits (MLP) & \Cell & \Cell & \Cell & \Cell & \Cell & \Cell & \Cell & \Cell \\
synthetic tabular (MLP) & \Cell & \Cell & \Cell & \Cell & \Cell & \Cell & \Cell & \Cell \\
breast-w & \Cell & \Cell & \Cell & \Cell & \Cell & \Cell & \Cell & \Cell \\
climate-model-simulation-crashes & \Cell & \Cell & \Cell & \Cell & \Cell & \Cell & \Cell & \Cell \\
credit-approval & \Cell & \Cell & \Cell & \Cell & \Cell & \Cell & \Cell & \Cell \\
ionosphere & \Cell & \Cell & \Cell & \Cell & \Cell & \Cell & \Cell & \Cell \\
sonar & \Cell & \Cell & \Cell & \Cell & \Cell & \Cell & \Cell & \Cell \\
wdbc (sensitivity node) & \Cell & \Cell & \Cell & \Cell & \Cell & \Cell & \Cell & \Cell \\
\midrule
Pooled & \Cell & \Cell & \Cell & \Cell & \Cell & \Cell & \Cell & \Cell \\
\bottomrule
\end{tabular}
\end{table*}

\begin{figure}[t]
\centering
\FigShell{false-keep rate per node against the node's noise level, both proposers}
\caption{Reading guide: each point is a node under one proposer; the horizontal axis is the comparison noise relative to the typical increment, and the vertical axis is the false-keep rate of the default loop.}
\label{fig:c1}
\end{figure}

\paragraph{False and missed keeps.}
Pooled over nodes, \COneFalseKeepSOne{} of the keeps in sweep~1 and \COneFalseKeepSTwo{} in sweep~2 had an audited increment that was not positive. Across nodes the rate ranges over \COneFalseKeepRange. Figure~\ref{fig:c1} plots it against each node's noise level. \TBD{One sentence on the relation in the figure, no stronger than eight nodes allow.} Our false-keep rate is a share of keeps, whereas \citet{jalaian2026noise} report the share of truly neutral candidates that a rule keeps, so the two are not directly comparable. The missed-keep rate is \COneMissedKeepSOne{} in sweep~1 and \COneMissedKeepSTwo{} in sweep~2.

\paragraph{Record bias.}
The record bias, the recorded gain of the incumbent minus its audited gain, is \COneRecordBiasSOne{} in sweep~1 and \COneRecordBiasSTwo{} in sweep~2, in units of $\sigma_{\mathrm{node}}$. Proposition~\ref{prop:bias} states that a keep leaves the record with a positive expected error whenever the candidate's error has conditional mean zero and a discard was possible. The measurement gives the size of the bias in the logged campaigns. It is not a test of the proposition, which is a mathematical statement. With a noise component shared by the campaign, the zero-mean condition can fail (Section~\ref{sec:theory-bias}). \TBD{Sign and size per node; name every node on which the bias is not positive.}

\paragraph{Wasted runs.}
Runs on duplicates, no-ops and invalid proposals take \COneWastedSOne{} of the runs in sweep~1 and \COneWastedSTwo{} in sweep~2. Duplicates alone take \COneDupSOne{} and \COneDupSTwo. Table~\ref{tab:app-dup} gives the duplicate share per node, sweep and memory mode. \TBD{Duplicate share with and without memory, per sweep.} \citet{chen2026visible} report duplicate shares in a deterministic setting (Section~\ref{sec:related}). Their loop, proposer and denominator differ from ours, so we give the numbers side by side and do not compare them formally. On a node with bit-exact runs, a duplicate returns the metric already in the ledger and cannot be kept under the default rule (Proposition~\ref{prop:app-rules}), so its run is lost whatever the noise. The rule layer of A2 skips such proposals; Sections~\ref{sec:results-gate} and~\ref{sec:results-closed} include its value offline and in the closed loop. On the exam data, the rule layer had less to skip for the main proposer than for the sweep-1 manager (Appendix~\ref{app:exam}). \TBD{Whether the two sweeps show the same ordering, descriptively.}

\paragraph{Controls.}
% Revised 23 Nov 2026 (R1-CONTROLS-NEG-DROP): the negative control is dropped; AI-assisted.
A positive control supports the reading of Table~\ref{tab:c1} (Table~\ref{tab:app-controls}). It asks whether the audit can see a known effect. For each node, it takes the single edits of the default configuration whose effects were measured with 30 seeds in a separate noise-floor run, and gives the probability that an audit with $K=10$ seeds, and with the $K=5$ seeds that the gate's labels use, read with the audit statistic that C1 uses, detects each edit. The probability is estimated by resampling 10 (or 5) of the 30 seeds. It is \COnePositiveControl. The control is not a test, and it reads no test label and no evaluation value. A negative control on edits that change nothing is not possible here: the harness rejects such an edit before running it, so neither sweep executed one, and on nodes whose runs are bit-exact the default rule cannot keep a repeated configuration (Proposition~\ref{prop:app-rules}).

% MOVED from results.tex (C3 and the secondary analyses). The paragraph on what carries the paper is rewritten.
\subsection{The offline ladder and the gate (C3)}
\label{app:gate-full}

The gate was read once, on \GateReadDate, after the arms-frozen record had fixed every arm, split and gate file \TBD{citations of the arms-frozen record and of the test release record}. The read covers the seven primary T0 nodes and \TOneNodesUsed. \GateEligibleNodes{} met the eligibility conditions of Section~\ref{sec:setup} and entered the IQM; \GateExcludedNodes. The Holm step-down ranked \GateHolmFirst{} first, so that arm was tested at the 97.5\% level and the other at 95\%. The primary test input leaves out test proposals that repeat a training or development proposal; this removed \GateDupLeftOut. Repeats concentrate on early proposals, because every instance of a node starts from the same configuration, so the primary estimand weighs later proposals more than the full test population does. \TBD{Counts of partial and divergent proposals on the evaluation seeds, per arm. If utilities tie: the number of ties, and the statement that the PoI counts a tie as one half and that the calibration evidence of the pass rule covers continuous utilities only.}

\begin{table*}[t]
\caption{The offline ladder on the test instances of the primary input. Reading guide: $U$ is the offline one-step utility scaled to 20 proposals, in units of $\sigma_{\mathrm{node}}$, averaged over nodes; gap closed and PoI are relative to A2; lower bounds are one-sided at the Holm level of the arm (97.5\% for the first-ranked arm, 95\% for the second). Only A3 and A5 carry a verdict, and every other row is descriptive. A6 is not in the arms-frozen record and has no offline test read. \TBD{Interval method for the descriptive rows; confirm that the test reads of A0, A1 and the secondary rows are inside the preregistered scope; final list of secondary rows as in Section~\ref{sec:setup}.}}
\label{tab:ladder}
\centering
\small
\begin{tabular}{llcccc}
\toprule
Arm & Forecast & $U$ & IQM gap closed [lower] & PoI [lower] & Gate \\
\midrule
A0 & none (default loop) & \Cell & \Cell & \Cell & not tested \\
A1 & constant prior & \Cell & \Cell & \Cell & not tested \\
A2 & A1 with rule layer & \Cell & reference & reference & comparator \\
A3 & tabular & \Cell & \Cell & \Cell & \Cell \\
A4 & proposer zero-shot & \Cell & \Cell & \Cell & not tested \\
A5 & fine-tuned small model & \Cell & \Cell & \Cell & \Cell \\
A6 & A5 + decision-focused fine-tuning & \multicolumn{4}{c}{no offline test read} \\
\midrule
\multicolumn{6}{l}{\emph{Secondary rows (descriptive)}} \\
A5, other size & fine-tuned small model & \Cell & \Cell & \Cell & not tested \\
A5, base model & fine-tuned base model & \Cell & \Cell & \Cell & not tested \\
TabPFN & tabular foundation model & \Cell & \Cell & \Cell & not tested \\
\midrule
Oracle & own label block & \Cell & 1 by definition & \Cell & not tested \\
\bottomrule
\end{tabular}
\end{table*}

\begin{figure}[t]
\centering
\FigShell{utility of each arm with bootstrap intervals, per tier}
\caption{Reading guide: arms are ordered as in Table~\ref{tab:ladder}; only A3 and A5 carry a verdict, and the intervals of the other arms are descriptive.}
\label{fig:ladder}
\end{figure}

Table~\ref{tab:ladder} gives every arm's utility and, for A3 and A5, the gate statistics with their lower bounds. Only A3 and A5 carry a verdict; every other row is descriptive. A6 has no offline values. It is trained only after a gate that passes, so it is not in the arms-frozen record, and a test read of it would need an amendment of the preregistration. The development skill that selected each learned arm's training source and, for A5, its model size is recorded in the arms-frozen record (Table~\ref{tab:app-devskill}).

\ifGatePassed
\paragraph{The gate passes.}
For A3, the IQM of gap closed is \GateIQMAThree{} and the PoI is \GatePoIAThree; for A5 they are \GateIQMAFive{} and \GatePoIAFive. Lower bounds are one-sided at each arm's Holm level. Under the Holm step-down, the pass rule is met by \GatePassingArms. The claim that the preregistration supports is limited to a passing arm. On the nodes of the read and at the declared run price, the arm's forecast, fed to the closed-form rule, earns more offline one-step utility than A2's constant forecast in the two senses that the pass rule tests, each through its lower bound: the interquartile mean of the gap closed over eligible nodes is above zero, and, averaged over nodes, the arm beats A2 on more than half of the test instances. An arm that does not pass is reported, not claimed. The point IQM of a passing arm is \TBD{above or below} the practical size of 0.1 declared with the gate. \TBD{If below: the gain is detectable but smaller than the declared practical size, and the text says so.} The gate has low power (Section~\ref{sec:setup}), and the estimate of a test that passes at low power tends to be too large, so we take the lower bound, not the point value, as the size that the data support.

\paragraph{The ladder.}
\TBD{One paragraph, descriptive: the step from A0 to A2; the step from A2 to the passing arm; A3 against A5, with no test, since their contrast is not preregistered; A4, the proposer's own forecast, against A5; the secondary rows. Name arms by their rows, give intervals, and rank no further than the intervals allow.}
\else
\paragraph{The gate does not pass.}
Neither A3 nor A5 passes. For A3, the IQM of gap closed is \GateIQMAThree{} and the PoI is \GatePoIAThree; for A5 they are \GateIQMAFive{} and \GatePoIAFive. Lower bounds are one-sided at each arm's Holm level. \TBD{Name the condition that failed for the first-ranked arm: the IQM bound, the PoI bound or both. Under the step-down, the second arm cannot pass once the first fails.} We report this as the measured negative of C3: on these nodes, neither tested arm showed a gain over A2 that the preregistered rule detects. This is the absence of a demonstrated gain, not evidence that a learned forecast cannot add utility.

\paragraph{Power.}
In the simulation that calibrated the pass rule, the mean two-arm power was \GatePowerMain{} at an effect size of \GatePowerEffect, and \GatePowerSmallTOne{} when a T1 node has six instances (Section~\ref{sec:setup}). The rule was chosen for its false-pass rate, and part of its calibration comes from the screen on positive denominators. In that simulation, a true gap closed of \GatePowerEffect{} with a headroom of at least 1 (Section~\ref{sec:setup}) fails the gate more often than it passes. \TBD{The realised configuration (eligible nodes per tier, test instances per node) beside the simulated one, without a recomputed power: the preregistration reports the simulated power only.}

\paragraph{Effect sizes.}
Table~\ref{tab:ladder} gives the effect size of every arm with its interval. \TBD{Interval method for the rows that the gate does not test, and whether upper bounds from the gate's bootstrap resamples may be reported; the preregistration fixes only the one-sided lower bounds of A3 and A5.} \TBD{If upper bounds are reported: the largest IQM of gap closed of A3 and of A5 that the data do not exclude.} We do not interpret differences among the learned arms. \TBD{Descriptive: the step from A0 to A2, and A4 against A5.}

\paragraph{The dispersion bound, as an estimate.}
Proposition~\ref{prop:app-gap} bounds the loss of a keep decision made with a constant threshold in place of the per-proposal threshold. The loss is quadratic in the distance between the two thresholds. A per-proposal forecast can therefore add much to the keep decision only where the per-proposal thresholds spread widely around the node's constant threshold. We evaluate the upper bound with the fitted forecasts of A3 and A5 in place of the true parameters, against A2's constant threshold, on \TBD{proposal population}. Scaled to 20 proposals and in units of $\sigma_{\mathrm{node}}$, it is \DispBoundAThree{} with A3's forecast and \DispBoundAFive{} with A5's (per node in Table~\ref{tab:app-dispersion}). We read this number as an estimate, not as a bound on what a better forecaster could add. The proposition bounds the loss in terms of the true parameters; with fitted parameters, the error of the estimate is not controlled. A forecast that is too smooth understates the spread of the per-proposal thresholds, and so understates the room for improvement. The bound covers the keep threshold only, at fixed execution decisions; screening, whose one-step gain is $(g-V^*_n)^+$ (Proposition~\ref{prop:onestep}), lies outside it. Its campaign-level reading needs (NF) and (NI), which the loops we run do not satisfy, so we read it per proposal. The estimate uses forecasts and the measured noise only and reads no label. It is descriptive and is not part of the preregistration. \TBD{Check status of Proposition~\ref{prop:app-gap} and of the value formula it rests on (Proposition~\ref{prop:app-threshold}): author-checked, AI-checked under the proof policy, or not checked.}

\paragraph{The positive control and the negative.}
% Rewritten 26 Nov 2026 (R1-ALIGN-OUTCOME): the 23 Nov paragraph ``What carries the paper'' said that the
% paper then rests on C1, C2 and the rule arms. The headline is now the closed loop, which includes the
% learned arms whatever the gate shows (Section sec:results-closed). The control sentences are kept.
A gate that does not pass leaves the one-step claim without a demonstrated gain. The closed loop is read on its own terms (Section~\ref{sec:results-closed}). The positive control of C1 bears on the reading of the negative. \TBD{Per node: if the audit detects the known single edits with high probability, the negative is not explained by audits that cannot resolve effects of that size.} The gate's labels and values each use five of the ten audit seeds, so the control bears on them only indirectly.
\fi

\subsection{Secondary analyses}
\label{sec:results-secondary}

Every analysis in this subsection is descriptive and carries no verdict. \ifGatePassed None of them adds to or narrows the claim of the primary read.\else None of them can turn the primary read into a pass.\fi

\begin{table}[t]
\caption{Sensitivity of the gate statistics. Reading guide: the first row is the primary read; every other row is descriptive and carries no verdict.}
\label{tab:sens}
\centering
\footnotesize
\setlength{\tabcolsep}{2pt}
\begin{tabular}{L{1.0in}cccc}
\toprule
 & \multicolumn{2}{c}{A3} & \multicolumn{2}{c}{A5} \\
Run & IQM & PoI & IQM & PoI \\
\midrule
Primary: seven T0 nodes and the T1 nodes, repeats left out & \Cell & \Cell & \Cell & \Cell \\
Eight T0 nodes (with \texttt{openml\_wdbc}) & \Cell & \Cell & \Cell & \Cell \\
Full test population (repeats kept) & \Cell & \Cell & \Cell & \Cell \\
Run price $0.05\,\sigma_{\mathrm{node}}$ & \Cell & \Cell & \Cell & \Cell \\
Run price $0.2\,\sigma_{\mathrm{node}}$ & \Cell & \Cell & \Cell & \Cell \\
Memory: append-only summary only & \Cell & \Cell & \Cell & \Cell \\
Memory: none only & \Cell & \Cell & \Cell & \Cell \\
\bottomrule
\end{tabular}
\end{table}

\paragraph{Sensitivity.}
Table~\ref{tab:sens} recomputes the gate statistics under the declared variations: the eight-node run that adds \texttt{openml\_wdbc}, the full test population with repeats kept, the two sensitivity run prices, and each memory mode alone. \TBD{One paragraph: whether adding \texttt{openml\_wdbc} or keeping repeats changes the reading; how the statistics move with the run price; whether the gain differs by memory mode.}

\begin{table}[t]
\caption{Leave-one-suite-out on T1. Reading guide: each row holds out one suite; fold arms are retrained without it; all entries are descriptive and carry no verdict. If a suite had no labelled T1 node when the arms were frozen, the analysis stops as a whole and this table is replaced by that statement.}
\label{tab:loso}
\centering
\small
\setlength{\tabcolsep}{3pt}
\begin{tabular}{lcccc}
\toprule
 & \multicolumn{2}{c}{A3} & \multicolumn{2}{c}{A5} \\
Held-out suite & IQM & PoI & IQM & PoI \\
\midrule
autoresearch & \Cell & \Cell & \Cell & \Cell \\
MLGym & \Cell & \Cell & \Cell & \Cell \\
MLAgentBench & \Cell & \Cell & \Cell & \Cell \\
\bottomrule
\end{tabular}
\end{table}

\paragraph{Transfer to a held-out suite or node.}
In leave-one-suite-out (Table~\ref{tab:loso}), the fold arms are retrained without the held-out suite and scored on its nodes. \TBD{If every suite has a labelled T1 node in the arms-frozen record and every fold arm is in that record: one descriptive paragraph on whether A3 and A5 keep their utility on a held-out suite. Otherwise delete the table and write that leave-one-suite-out stopped as a whole, because a named suite had no labelled T1 node when the arms were frozen, or because a fold arm was not in the arms-frozen record, and that the aggregate gate and the other secondary analyses are unaffected.} Leave-one-node-out within T0 gives each held-out node's gap closed, with fold arms retrained without that node (Table~\ref{tab:app-lono}).

\paragraph{Secondary arms and the second proposer.}
The secondary arms are compared with A2 per node (Table~\ref{tab:app-secondary}): A4, the other model size of A5, the same fine-tuning applied to the base model, and the tabular foundation model \citep{hollmann2025tabpfn}. A4 asks how far the proposer can judge its own proposals; we report it beside A5 without a test. \TBD{A4 against A5, descriptive.} Sweep~1 enters the analysis only as a descriptive stratum (Table~\ref{tab:app-sweep1}). A table that transfers arms from one proposer's training instances to the other proposer's test instances is outside the preregistered scope and was not read. The per-node gate table, with the numbers of training, development and test instances, the eligibility statistics and the denominators, is Table~\ref{tab:app-gate-nodes}.
